%% Citas verificadas contra el PDF. Sin marcadores [VERIFY] pendientes.

\section{TRABAJO RELACIONADO}

La estimación de riesgo en comercio electrónico constituye un campo de investigación
consolidado. Esta sección lo analiza en cinco dimensiones: inteligencia artificial para
estimación de riesgo, detección de vendedores fraudulentos desde la plataforma,
identificación de reseñas manipuladas, detección de sitios de suplantación y contramedidas
implementadas por las organizaciones comerciales. Un apartado final identifica la brecha que
este trabajo aborda. Los avances se concentran en el lado del operador de la plataforma, y
rara vez llegan al comprador en el momento de la decisión de compra.

\subsection{Inteligencia artificial aplicada a la estimación de riesgo}

El aprendizaje profundo se ha consolidado como la tecnología preferida para tareas como
estimación de riesgo y toma de decisiones en plataformas de comercio electrónico \cite{ref31}.
Los sistemas basados en reglas no escalan ante el volumen y el carácter adaptativo de las
amenazas actuales. El contenido generado por usuarios como reseñas, calificaciones e
imágenes resulta decisivo en la decisión de compra y, simultáneamente, altamente vulnerable a
la manipulación y al fraude coordinado.

Un ejemplo cuantificado ilustra tanto el nivel alcanzado como las condiciones que lo
hacen posible. Un modelo híbrido de red convolucional con máquina de soporte vectorial
optimizada alcanza una exactitud del 96.23\,\% y un AUC de 0.9118 al clasificar el riesgo
crediticio de empresas de comercio transfronterizo \cite{ref32}. Sin embargo, el objeto
estimado no es el nivel de riesgo de estafa que percibe un comprador individual. Los insumos
que ese desempeño requiere información financiera de la contraparte, condiciones de
operación y evaluaciones de la plataforma no son accesibles desde el navegador de quien está
por realizar una compra \cite{ref32}.

Los algoritmos que sustentan esta línea de investigación son identificables. En la detección
de reseñas manipuladas predominan los modelos de lenguaje basados en transformadores:
variantes de BERT ajustadas al dominio del comercio electrónico \cite{ref16}, arquitecturas
adversariales sobre la misma familia \cite{ref17} y redes neuronales profundas aplicadas a
reseñas de plataformas concretas \cite{ref18}. En ocasiones se combinan con memorias de largo
y corto plazo bidireccionales y con herramientas de explicabilidad para justificar el
veredicto \cite{ref16}. En la estimación de riesgo del vendedor y en la detección de sitios de
suplantación predominan, en cambio, los métodos de ensamble y los clasificadores de margen:
bosques aleatorios, máquinas de soporte vectorial, árboles de decisión y gradiente potenciado
aparecen de forma recurrente \cite{ref10, ref13, ref14}. El marco de detección más reciente
combina cuatro de ellos con selección de características mediante SHAP y eliminación recursiva
\cite{ref25}. Esta distribución no es casual: el texto libre de las reseñas exige
representaciones contextuales, mientras que los atributos estructurales de una ficha
heterogéneos y de baja dimensionalidad se ajustan mejor a los métodos basados en árboles.

Estos enfoques se sitúan invariablemente en la plataforma, en la institución financiera o en
el operador del servicio, que son quienes disponen de los registros y ejecutan los modelos.
Incluso en esos entornos admiten limitaciones: los detectores más avanzados pueden ser eludidos
mediante perturbaciones adversariales u ofuscación lingüística \cite{ref31}. La
defensa centralizada, por sólida que sea, deja pasar casos que alcanzan al comprador sin que
este disponga de instrumento alguno de protección.

\subsection{Detección de fraude desde la plataforma}

La detección de fraude en comercio electrónico ha alcanzado madurez metodológica \cite{ref10}.
Un exponente reciente combina razonamiento basado en casos con un sistema multiagente para
alcanzar un AUC de 0.979 sobre 1.27 millones de instancias, resistiendo la deriva del
comportamiento fraudulento \cite{ref28}. Dicho desempeño se sostiene sobre datos que solo el
operador posee: registros transaccionales, huellas de dispositivo y un grafo de interacción de
el comprador que evalúa una ficha no tiene acceso.

\subsection{Reseñas manipuladas}

La segunda línea de investigación aborda la autenticidad de las reseñas. Además de los
detectores basados en contenido textual, esta línea incorpora enfoques de fusión que combinan
diversas características de la reseña \cite{ref15}. Un modelo de competencia demuestra que en
equilibrio el vendedor auténtico no adquiere reseñas falsas mientras que el falsificador puede
hacerlo para inducir a error al comprador \cite{ref29}. La manipulación de la reputación no es,
por tanto, una anomalía sino un comportamiento previsible. La cantidad de reseñas falsas en
equilibrio decrece con la fracción de consumidores capaces de interpretar correctamente los
indicadores de la ficha. Esto sitúa la información del comprador y no
solo la capacidad de detección de la plataforma como palanca de mitigación efectiva.

\subsection{Suplantación de sitios y análisis de dominio}

La tercera línea proviene de la seguridad informática. En \cite{ref7} se detectan localizadores
uniformes de recursos (URL) de \textit{phishing} en dominios aparcados a partir de
características derivadas del registro, aplicándolo a una institución financiera; en
\cite{ref25} se propone un marco de detección sobre 1\,353 URL con selección de
características por explicaciones aditivas de Shapley (SHAP) y eliminación recursiva, alcanzando
una exactitud del 98.4\,\%.

Lo más relevante de este último para el presente trabajo no es el desempeño del clasificador
sino el diagnóstico que lo motiva: el certificado SSL válido dejó de discriminar sitios
legítimos de fraudulentos. Los atacantes obtienen certificados mediante servicios automatizados,
y las listas de bloqueo del navegador no alcanzan a cubrir dominios desplegados y retirados en
horas \cite{ref25}.

\subsection{Contramedidas del lado del comercio}

Una revisión sistemática de métodos de seguridad en comercio electrónico \cite{ref30} inventaría
siete familias de controles y concluye que ninguna basta por sí sola, por lo que la
defensa debe articularse en capas. Las siete, no obstante, se implementan y financian dentro
de la organización comercial. Ninguna protege al comprador que interactúa con una plataforma o
un vendedor que no las ha adoptado.

\subsection{Brecha de investigación}

El comprador ocupa una posición singular: existen métodos eficaces contra vendedores
fraudulentos, reseñas manipuladas y suplantación de sitios, pero todos requieren datos,
infraestructura o autoridad de los que él carece. Las recomendaciones dirigidas a él son de
naturaleza educativa y de eficacia limitada \cite{ref27, ref3}. Ninguno de los trabajos
revisados combina características de dominio, de vendedor y de reseñas en un instrumento que
se ejecute del lado del comprador y emita su advertencia antes del pago. Esa es la brecha que
este trabajo aborda. Las características propuestas son trasladables a cualquier plataforma de
comercio electrónico que exponga fichas de producto con información pública: MercadoLibre,
Amazon, eBay o AliExpress, entre otras.
